\chapter{绪论}
\setlength{\parskip}{0.4em}
\setlength{\itemsep}{0.2em}
\setlength{\tabcolsep}{4pt}

\section{数字信号与数字电路}
\subsection{模拟量与数字量}
自然界的物理量就其变化规律而言不外乎两大类。

\begin{itemize}
    \item \textbf{模拟量}：在时间上或数量上的变化是连续的。如加热炉内的温度、水库水位的高度，相邻两值之间总可插入无穷多个中间值。
    \item \textbf{数字量}：在时间上和数量上都是离散的，变化发生在不连续的瞬间，且数值的大小与每次的增减都是某个\textbf{最小数量单位} $\Delta$ 的整数倍。如统计每天从装配线上输出的汽车数量。
\end{itemize}

把模拟量、数字量转换成电压（或电流）信号，得到的信号也分为两大类：\textbf{模拟信号}表征模拟量，幅值随时间连续变化；\textbf{数字信号}表征数字量，只取有限个离散状态，在二进制数字电路中仅取高、低两种电平，用 $1$、$0$ 表示。工作在模拟信号下的电子电路称为\textbf{模拟电路}，工作在数字信号下的电子电路称为\textbf{数字电路}。

需要强调：数字信号中的 ``$1$'' 和 ``$0$'' \textbf{没有数量的意义}，只表示事物的两个对立面（高/低电平、开关通/断、灯亮/灭），其数量含义要由\textbf{数制}和\textbf{码制}另行赋予。

\subsection{数字电路与模拟电路的比较}
二者在工作信号、研究对象、分析设计方法以及所用的数学工具上都有明显不同，见表~\ref{tab:dig_vs_ana}。

\begin{table}[!htbp]
    \centering
    \caption{数字电路与模拟电路的比较}
    \label{tab:dig_vs_ana}
    \begin{tabular}{@{}lll@{}}
        \toprule
        对比维度 & 数字电路 & 模拟电路 \\
        \midrule
        工作信号 & 离散的高、低电平序列 & 连续的电压、电流波形 \\
        取值特点 & 只有 $0$、$1$ 两种状态，可容偏差 & 连续取值，精度依赖器件一致性 \\
        研究对象 & 输入输出的逻辑关系与时序关系 & 信号的放大、变换与波形质量 \\
        分析方法 & 逻辑代数、真值表、卡诺图、状态转换 & KCL/KVL、小信号等效电路、频响分析 \\
        数学工具 & 离散数学（布尔代数） & 微积分、复变函数、线性代数 \\
        \bottomrule
    \end{tabular}
\end{table}

\subsection{数字电路的主要优点}
\begin{enumerate}
    \item \textbf{抗干扰、容偏差}：每一位只有 $0$、$1$ 两种取值，电路只要能正确区分两种状态即可，无须精确复制某个电压值，因而对制造精度、工作条件和运行环境的要求大大降低。
    \item \textbf{精度可扩展}：需要提高精度时只需增加二进制数的位数，不必提高单个器件的精度。
    \item \textbf{便于集成与存储}：单元电路重复性强，易于大规模集成；信息以代码形式保存，便于长期存储与复制。
\end{enumerate}
正因为对制造精度要求相对宽松，\textbf{首先制成的集成电路是数字集成电路}；迄今为止大多数大规模、超大规模集成电路都属于数字集成电路。

为充分发挥数字电路在信号处理上的优势，可先将模拟信号按比例转换成数字信号，送入数字电路（或计算机）处理，再把结果转换成模拟信号输出。自 20 世纪 70 年代以来，这种用数字电路处理模拟信号的``数字化''浪潮已席卷几乎所有电子技术应用领域。

\section{电子技术的发展历程}
电子技术的发展与进步，是与电子器件的更新紧密相连的。表~\ref{tab:dev_history} 按代表性器件给出四个阶段的划分。

\begin{table}[!htbp]
    \centering
    \caption{电子技术发展的四个阶段}
    \label{tab:dev_history}
    \begin{tabular}{@{}llll@{}}
        \toprule
        阶段 & 年代 & 代表器件 & 主要特征与应用 \\
        \midrule
        第一阶段 & 20 世纪 40 年代 & 电子管（真空管） & 早期计算机、电话交换、数字通信 \\
        第二阶段 & 20 世纪 60 年代 & 晶体管（半导体三极管） & 功耗小、寿命长、体积小，扩展到测量领域 \\
        第三阶段 & 20 世纪 70 年代中期 & 集成电路（IC） & 数字技术广泛应用，进入医疗、雷达、卫星 \\
        第四阶段 & 70 年代中至 80 年代中 & 大规模、超大规模 IC & 微电子技术迅猛发展，渗透到各行各业 \\
        \bottomrule
    \end{tabular}
\end{table}

\subsection{关键器件与事件}
\begin{itemize}
    \item \textbf{1904 年，真空二极管}：19 世纪 80 年代爱迪生发现``爱迪生效应''，弗莱明据此发明真空二极管；德福雷斯特在灯丝与板极之间加入栅极，发明第一只真空三极管，被看作电子工业真正的诞生起点。电子管工作时须将阴极加热到很高温度，因而功耗大、寿命短、体积重量大、噪声大、工艺复杂。
    \item \textbf{1947 年，晶体管}：布拉顿（W. Brattain）、巴丁（J. Bardeen）和肖克利（W. Shockley）发明半导体三极管，功耗小、寿命长、体积重量大大减小，几乎在所有领域逐步取代电子管，发明者获诺贝尔奖。
    \item \textbf{1958 年，集成电路}：基尔比（J. Kilby）将若干晶体管、电阻及连线制作在一小片硅片上；同一时期诺伊斯（R. Noyce，Fairchild）也独立实现了集成电路，开辟了电子电路微型化的新途径。
    \item \textbf{1965 年，摩尔定律}：摩尔（G. E. Moore）提出集成电路中晶体管数目将以每年翻一番的速度增长（20 世纪 80 年代末以后翻番时间放慢到约 18 个月），同时每个晶体管的价格降为原来的一半。这一预测被实践所证明，称为``摩尔定律''。
    \item \textbf{当代}：已能把上亿个晶体管集成在一片邮票大小的硅片上，甚至把整个系统集成为``片上系统''（System on Chip，SoC）。2010 年集成电路微细加工精度已达 32\,nm；集成度的提高最终受加工工艺极限与晶体管物理极限的限制，摩尔定律也终将失效。
\end{itemize}

\subsection{电子电路的核心功能}
电子电路的功能可统一概括为\textbf{信息处理}与\textbf{能量转换}：模拟电路用连续的模拟电压、电流值表示信息；数字电路用一个离散的电压序列（$0/1$ 代码序列）表示信息。

\section{数字电路的特点与分类}
\subsection{数字电路的特点}
\begin{enumerate}
    \item \textbf{开关工作}：半导体器件工作在开关状态（截止或饱和/导通），只区分高、低两种电平。
    \item \textbf{二值描述}：输入、输出量都可用二值逻辑变量表示，输入与输出之间是逻辑上的因果关系，可用逻辑函数 $Y=F(A,B,C,\cdots)$ 描述。
    \item \textbf{研究重点}：不追求波形保真，而关心逻辑关系是否正确、时序是否满足要求、电平是否符合接口规范。
\end{enumerate}

\subsection{分类}
\begin{itemize}
    \item 按处理的信号形式：模拟电路、数字电路。
    \item 按输出与历史状态是否有关：\textbf{组合逻辑电路}（输出仅取决于该时刻输入）与\textbf{时序逻辑电路}（输出还与电路原状态有关）。
    \item 按所用器件：\textbf{双极型}（TTL、ECL 等）与\textbf{单极型}（CMOS 等）。
    \item 按集成度：\textbf{小规模}（SSI，基本门、触发器）、\textbf{中规模}（MSI，译码器、数据选择器、计数器）、\textbf{大规模}（LSI，存储器、可编程逻辑器件）、\textbf{超大规模}（VLSI，CPU、SoC）。
    \item 按层级：单体 MOSFET $\rightarrow$ 逻辑门 $\rightarrow$ 组合/时序功能块 $\rightarrow$ 复杂数字系统。
\end{itemize}

\section{可编程逻辑器件与 EDA 技术}
\subsection{从固定逻辑到可编程逻辑}
早期数字集成电路的逻辑功能在制造时即已固定，要改变功能就必须改变内部单元电路之间的连接，而连接在制作过程中已经固定下来。若制作一种允许用户自行修改内部连接的通用集成电路，则只需生产这一种器件就能满足不同用户的需要，这种器件称为\textbf{可编程逻辑器件}（Programmable Logic Device，\textbf{PLD}）。PLD 的内部电路结构可通过写入编程数据来设置，部分器件的编程数据还可擦除重写。PLD 最初出现在 20 世纪 60 年代后期，80 年代以后获得迅猛发展。

需要注意：微处理器和单片机也是可编程的大规模数字集成电路，但对它们编程并不改变电路的内部连接结构，只是写入运行程序，其内容属于``微机原理''等后续课程。

\subsection{EDA 技术}
如此复杂的编程工作无法用手工完成，PLD 的编程（包括编程数据的生成和写入）必须使用 \textbf{EDA}（Electronics Design Automation，电子设计自动化）技术。其工具包含硬件和软件两部分：
\begin{itemize}
    \item \textbf{硬件}：计算机和编程器；
    \item \textbf{软件}：用于描述电路功能的\textbf{硬件描述语言}（Hardware Description Language，\textbf{HDL}）及其编译程序，以及适应于选用器件的编程软件，用以完成编程数据的生成与写入。
\end{itemize}

早期 PLD 需先把计算机生成的编程数据输入编程器，再写入置于编程器上的器件；如今大规模集成的 PLD 越来越多地采用``\textbf{在系统可编程}''（In System Programmable，\textbf{ISP}）技术，编程数据可直接写入已安装在系统里的器件，无须专用编程器，也不必把器件从系统中取出。

本课程涉及的现代设计流程可概括为：逻辑功能描述（HDL）$\rightarrow$ 逻辑综合 $\rightarrow$ 仿真验证 $\rightarrow$ 布局布线 $\rightarrow$ 下载/编程 $\rightarrow$ 硬件测试。

\section{课程内容与学习方法}
\subsection{课程体系中的位置}
先修\textbf{电路原理}；平行\textbf{模拟电子技术基础}及其实验、\textbf{数字电子技术基础}及其实验；后续\textbf{微机原理}、\textbf{嵌入式系统}。

\subsection{课程目的与基本要求}
课程目的是掌握\textbf{基本概念、基本分析方法、基本设计方法和基本实验技能}，可分解为四项核心能力：
\begin{enumerate}
    \item 基础概念：逻辑电平、逻辑门、组合/时序、噪声容限等；
    \item 电路分析方法：由电路求逻辑功能、由波形求真值表、时序分析；
    \item 电路设计方法：由功能要求选器件、化简、画逻辑图、验证；
    \item 实验技能：仪器使用、电路搭接与调试、故障定位。
\end{enumerate}

\subsection{学习方法}
\begin{itemize}
    \item \textbf{自底向上（Bottom-up）}：先理解器件与门如何工作，再逐级抽象到功能块与系统。
    \item \textbf{抓抽象与接口}：抽象方法与接口约定是复杂系统可管理的关键；接口应具备良好的隔离特性以及扩展、升级特性。
    \item \textbf{分层训练}：在各层级上掌握设计电路的方法，兼顾学科历史、系统方法、算法与电路调试方法。
    \item \textbf{重实践}：逻辑代数与化简必须动手推导，门电路与时序电路必须结合波形和实验现象理解。
\end{itemize}

优质数字系统设计的准则：结构化设计，每一级保持有限的复杂度，支持模块复用；以最小的代价实现最大的功能，实现可靠的系统，并保持对未来技术的兼容性。
